% !TeX root = ../papers/nc-p3.tex
\section{二次剰余の全放物線と経路行列の有理関数体への降下}
\label{sec:rank-fifteen-quadratic-parabolic-pencils}
\begin{theorem}
剰余標数3の付値で $X=\bar x\ne0,1$, $Y=\bar y\ne0$ とする。
$R=Y/(1-X)^2\in\mathbb F_9^\times$ なら、全ての通常な標数零の階数15窓は標準Koszul左右AS正則で、
各始点が四中心四次式上自由階数256である。
\end{theorem}
\begin{proof}
残る倍率の三乗共役の代表は $R=s+i$, $s=0,1,-1$, $i^2=-1$ である。
$r^2=R$, $a^3=r$, $\delta^2=R-1$ の全定数は $\mathbb F_{81}$ に取れる。
前節の $b=(q-q^3)/(r-a)$, $v=q+ab$, $Z=\delta(b^3-b)$ の表示で、
$B=-(q^3-q)^2=\kappa/X$, $\kappa=-(r-a)^2\in\mathbb F_9^\times$ である。
実際 $r^9=\pm r$ で $a^9$ にも同じ符号が付く。従って $B\in\mathbb F_9$ の有限非零水準は既知である。

\path{computations/verify_rank_fifteen_quadratic_parabolic_pencils.py} は
$\mathbb F_{81}(q)$ 上で五次までの経路と四次中心を計算する。
文字0の五本の自由列のうち第1,2,4,5列の中心を課すと、各始点の商が十三次まで標準次元を持つ。
下表で第3列は始点別の除外多項式の最小公倍式の次数、第4列は
$B\in\mathbb F_9$ を除いて $q\mapsto\pm q+j$ の六変換で最大公約元を取った次数、
第5列はその係数の三乗共役全体の和集合の次数である。
\[
\begin{array}{c|c|r|r|r|c}
R&\text{四始点の除外次数}&\text{最小公倍式}&\text{六変換}&\text{共役和}&\deg_B\\\hline
i &(1284,766,971,1142)&3831&66&84&4,4,6\\
1+i &(1476,743,1755,1723)&5416&60&96&4,4,8\\
-1+i &(1426,1136,1640,1844)&5768&60&72&4,4,4
\end{array}
\]
最終列は共役和を完全に覆う $B$ の素体上の既約因子の次数である。
各因子の根を九乗軌道に分け、$X=\kappa/B$, $Y=R(1-X)^2$ で元の七倍率を取り直す。
全軌道で四始点の十三次消滅までを直接確認する。最大の八次因子には
$\mathbb F_{3^{32}}$ を使う。六表示による被覆、係数の共役と標数零への持上げは
第\ref{sec:rank-fifteen-quadratic-verticals}節の証明と同じである。
\end{proof}

\begin{corollary}
3上に剰余次数2以下の素点を持つ数体 $L$ に対し、任意の $\lambda\in L^\times$ で
$y=\lambda(1-x)^2$ 上の非零幾何水準の通常窓は全て上の正則性と自由性を持つ。
さらに、通常な未認証点が残るなら $x,\lambda=y/(1-x)^2$ はともに代数的数であり、
各々のモニック最小多項式は3整で、定数項は3進単元、3での還元は $T^9-T$ と互いに素である。
\end{corollary}
\begin{proof}
前節の有理倍率の系と同じく、$\lambda$ が単元のときだけ今回の定理が必要である。
もし $x$ が超越的なら $3x$ に超越的剰余を与えるGauss付値を取り $\nu(x)<0$ とできる。
$\lambda$ が超越的なら同じく $\nu(\lambda)<0$ とでき、$\bar x\ne0,1$ の単元域で
$y$ は負付値となる。いずれも係数体への延長で全点を覆う。
代数的な $x,\lambda$ のいずれかが3上のある付値で非単元でも既知領域に入る。
全共役が単元ならモニック最小多項式は3整で定数項は単元である。
その還元に $\mathbb F_9$ の根があれば、代数閉包でモニック因子を一次式まで分解して還元すると、
その剰余を持つ共役がある。対応する付値で二次剰余の縦切片または本節の定理を使える。
\end{proof}

\begin{proposition}\label{prop:rank-fifteen-word-parity-descent}
36成分の正規化で、$F$ を六根の二乗と基底の三係数を含む体、
$E=F(W)$ を六根の積の添加とする。$W\ne0$ とする。任意の長さと始点について
文字別経路の関係行列、文字0の四次中心方程式、およびそこから選ぶ四中心商の関係行列は
語の基底変更と行のスカラー倍で $F$ 上に書ける。
特に二変数の通常表示では平方根の積を添加せず $F=k(K,H)$ でこれらの階数を計算できる。
\end{proposition}
\begin{proof}
$\omega(g,h)=g_L h_H\in\{0,1\}$ とし、位置 $p$ からの語 $w$ に
\[
 \ell_p(w)=\sum_j[\,p+j\equiv2\pmod4,\ w_j=3\,]+\sum_{j<k}\omega(w_j,w_k)\pmod2
\]
を付ける。添字 $j$ は0から始める。コードは全36係数の六根の共通指数の偶奇が
対応する $\ell_0(w)$ に等しいことを四つの巡回表示でも直接確認する。
連結 $UV$ の偶奇は両語の偶奇の和に $\omega(|U|,|V|)$ を加えたもので、$|U|$ はKlein次数である。

語基底を $W^{\ell_p(w)}w$ に変える。二次関係の固定した行は四語テンソルの残りの二文字が固定で、
変化する二文字のKlein次数は一定である。前後に固定した語を付けた任意の関係行でも、
各係数の $W$ 指数の偶奇はこの基底で行全体に共通になる。行をその偶奇の $W$ 冪で割ればよい。
従って各経路空間はこの基底で $F$ 上の商として計算できる。
四次の文字0の中心候補を同じ規格化語の線型結合とすると、矢を左右から掛けたときの
交差項 $\omega(0,h),\omega(h,0)$ はともに零である。
四位置周期性により矢の寄与も一致するので、交換子の行列も同じ操作で $F$ に降りる。
$F$ 上で解いた核は $E$ への拡大で全核を与える。得た文字0の切断の前後に語を
付けた行にも同じ議論が適用できる。$W^2\in F$ であることから結論が従う。
この議論は各長さでの線型代数を有理化する。積を再帰的に計算する場合、
次数 $g,h$ の積から $W^{\omega(g,h)}$ を除くと、括弧を変えたときの倍率は
\[
 (W^2)^{(\omega(h,k)+\omega(g,h+k)-\omega(g,h)-\omega(g+h,k))/2}.
\]
分子は偶数で、一つの次数が0なら零である。
\end{proof}
